\section{Task State, Trees, and Local Laws}\label{sec:min-framework}
\label{sec:min-manifesto-state}

This section defines the object C-TreePO learns and audits. The central pair
is a readout $f$ and a state map $g$. The readout $f$ is the task-facing
function: it returns the quantity the researcher wants to measure. The state
map $g$ is the compression operator: it maps a raw object, or an already
compressed state, into a smaller object that should be sufficient for $f$.
The goal is not lossless compression. The goal is oracle-preserving
compression:
\[
  f(g(x)) \approx f(x).
\]
When the input is split into parts, the stronger requirement is that the same
task value is preserved along the merge route:
\[
  f\!\left(g(g(u)\concat g(v))\right)
  \approx
  f(u\concat v).
\]

\subsection{Classical Warm-Up: Count-Min as a Mergeable Sketch}

Count-Min is the most compact example of the interface
\citep{CormodeMuthukrishnan2005}. Let $A$ be a multiset or stream segment of
items from an item universe $\mathcal I$. Fix $d$ hash functions
$h_j:\mathcal I\to\{1,\ldots,w\}$. The state map
$g$ sends $A$ to a counter matrix $C_A\in\mathbb N^{d\times w}$:
\[
  C_A[j,\ell]
  =
  \sum_{a\in A}\One\{h_j(a)=\ell\}.
\]
Equivalently, each observed item increments one counter in every row.

The merge rule is ordinary matrix addition. For disjoint stream segments
$A$ and $B$,
\[
  g(A\uplus B)=C_A+C_B.
\]
This equality is exact. No item order, document boundary, or original stream
record is needed once the counter state has been formed. The state is a
homomorphic image of the multiset under addition.

For a query item $q$, the readout is
\[
  f_q(g(A))
  =
  \widehat n_q(A)
  =
  \min_{j\le d} C_A[j,h_j(q)].
\]
If no other item collides with $q$ in a row, that row's counter equals the
true frequency $n_q(A)$. Collisions can only increase the counter, so the row
minimum controls the overestimate. Count-Min therefore separates two facts:
the state merge is exact, while the readout is approximate because the state
is compressed.

This is the pattern C-TreePO generalizes. Count-Min fixes a query family
$\{f_q:q\in\mathcal I\}$ and hand-designs a sufficient state $g$ for that
family. The designer knows before deployment that a counter matrix is the
right state. The correctness argument is not that $g(A)$ preserves the stream;
it plainly does not. The correctness argument is that $g(A)$ preserves enough
information for the frequency readouts after arbitrary parallel merges.

\subsection{From Analytic State to Learned State}

C-TreePO uses the same three slots but fills them differently.
\[
  \text{raw object}
  \xrightarrow{\;g\;}
  \text{state}
  \xrightarrow{\;\text{merge by }g\;}
  \text{root state}
  \xrightarrow{\;f\;}
  \text{task value}.
\]
In a classical sketch, the state is an engineered tuple, bitmap, register
array, counter matrix, or quantile list. In C-TreePO, the state can be a text
summary, an embedding sequence, a tensor field, or any object in the function
class chosen for the application. What makes it a sketch is not its type. What
makes it a sketch is the sufficiency contract with the readout and the fact
that states, not scalar answers alone, are merged.

This distinction is important. Scalar task values usually do not merge. The
frequency of $q$ happens to merge by addition, but many useful readouts do
not. Cardinality does not follow from $|A|$ and $|B|$ without knowing overlap.
A manifesto's economic-policy score does not follow from separately scoring
two halves and averaging them, because a fiscal qualification in one half can
change how a public-services pledge in the other half should be read. A
mergeable sketch preserves a richer state precisely because the scalar answer
is often underdetermined by child answers.

The hand-crafted-to-learned transition can be summarized as follows:
\[
  \begin{array}{lll}
  \text{Classical sketch:}
    & f \text{ fixed from a small query catalog,}
    & g \text{ analytically engineered for } f,\\[2pt]
  \text{C-TreePO:}
    & f \text{ supplied by an oracle or learned readout,}
    & g \text{ learned to be sufficient for } f.
  \end{array}
\]
The local laws are the replacement for the algebraic proof that a classical
sketch designer would normally write down.

\subsection{Learning $f$ and $g$}

Let $f^\star$ denote the trusted target available through expert labels,
rubric coders, a high-grade judge, or another costly oracle. In the ideal
case, the readout $f$ is fixed to this oracle. In practice, $f$ is often an
approximator trained or calibrated toward $f^\star$, while $g$ is trained to
produce states that are sufficient for $f$. A useful training objective has
two parts:
\[
  L(f,g)
  =
  L_{\mathrm{oracle}}(f,f^\star)
  +
  \lambda L_{\mathrm{local}}(f,g).
\]
The first term keeps the readout aligned with the target. The second term
penalizes violations of the local laws below. The tuning parameter $\lambda$
therefore controls the tradeoff between oracle fidelity and sufficiency of
the learned state. When $\lambda$ is too small, the readout may fit root labels
while the compression path remains unstable. When $\lambda$ is too large, the
system may prefer an easier-to-preserve surrogate target. The audit estimates
where the deployed pair landed.

For LLM-shaped states, learning can mean prompt search, DSPy/GEPA-style
program optimization, instruction tuning, supervised fine-tuning,
preference-based tuning, or other weight optimization against oracle and
local-law labels. The output $g(x)$ is usually text: a compact description of
the evidence the readout should need. This makes the state observable. The
researcher can inspect whether a leaf summary contains the load-bearing claim,
whether a merge summary preserved the interaction between two child summaries,
and whether two spans mapped to the same apparent state really are equivalent
for $f$.

For non-language domains, the same mathematical interface may be implemented
by a learned operator. If inputs are functions, fields, long trajectories, or
other structured objects, $g$ can be a neural operator mapping one function
space to another, and $f$ can be a learned readout on the compressed state
\citep{Kovachki23NeuralOperator,Li21FNO}. In that case the state may not be
human-readable text, but the local-law contract is unchanged: the learned
operator is useful only insofar as it preserves the oracle along the realized
merge route.

\begin{table}[t]
\centering
\small
\renewcommand{\arraystretch}{1.13}
\begin{tabular}{L{0.18\linewidth}L{0.24\linewidth}L{0.24\linewidth}L{0.23\linewidth}}
\hline
Regime & Readout $f$ & State map $g$ & Certificate \\
\hline
Count-Min & Fixed item-frequency query $f_q$ & Hand-designed counter matrix & Algebraic merge proof plus collision bound \\
Classical mergeable sketches & Fixed catalog query: cardinality, quantile, membership, moments & Engineered registers, tuples, bitmaps, or counters & Associative state merge and query-specific proof \\
LLM C-TreePO & Rubric scorer, calibrated judge, or learned readout toward an oracle & Prompted or weight-tuned readable text state & C1/C2/C3 audit, direct inspection, and design-based distortion estimate \\
General learned operator & Learned or calibrated task readout & Neural/functional operator between state spaces & Local-law residual on sampled realized calls \\
\hline
\end{tabular}
\caption{\textbf{C-TreePO generalizes mergeable sketches by learning the
state.} Classical sketches hand-design $g$ for a fixed $f$. C-TreePO learns
$g$ in a task-appropriate function class and audits whether the learned state
is sufficient for the readout.}
\label{tab:v5-learned-sketch-lineage}
\end{table}

\subsection{Running Example: Manifesto Policy State}

The running example is a manifesto's economic-policy position on the
public-services versus taxation scale used in the Manifesto/Benoit benchmark
\citep{BenoitEtAl2025,ManifestoProject2025a}. In this application, $f$ is the
target readout for that policy dimension. Given a span $x$, $f(x)$ is the
economic-policy score a trusted coder or calibrated rubric scorer would assign
after reading the relevant commitments to public services, taxation, welfare
expansion or retrenchment, market intervention, privatization, fiscal
restraint, and the qualifications that determine how those statements should
be interpreted.

The state map $g$ produces an observable policy-evidence state. A useful
$g(x)$ may discard slogans, duplicate rhetoric, document furniture, and
stylistic variation. It may not discard facts that can move $f$: whether a
tax cut is targeted or broad, whether a welfare commitment is universal or
means-tested, whether fiscal restraint is framed as temporary or permanent,
or whether market intervention is endorsed as a principle or only as an
emergency exception. Two spans can be different strings and still share a
task-relevant state if they support the same economic readout. Conversely,
two spans can look similar and fall in different fibers if one contains a
qualification the other omits.

The analogy to Count-Min is exact at the interface level. Count-Min stores a
counter matrix because counters are sufficient for frequency readouts.
C-TreePO stores a policy-evidence summary because that summary is intended to
be sufficient for the policy readout. Count-Min merges by adding matrices.
C-TreePO merges by applying $g$ to concatenated child states. Count-Min's
failure mode is collision error. C-TreePO's failure modes are missing
evidence, summary drift, and lost cross-span interaction.

\subsection{Compression Trees}

Let $\Strings$ be finite token sequences with concatenation
$u\concat v$. A summarizer is a possibly stochastic map
$g:\Strings\rightsquigarrow\Strings$; a readout
$f:\Strings\rightsquigarrow\Y$ scores the root state. In the Manifesto
application, $g$ is an LLM summarizer used at both leaves and merges, and $f$
is a rubric-facing scorer returning a policy score. In controlled examples,
$g$ may be a hand-designed sketch state, a learned tensor operator, or another
state map. The local laws attach to the pair $(f,g)$ and not to a particular
architecture.

A document $x$ is partitioned into contiguous leaves
$(b_1,\ldots,b_k)$ attached to a binary tree. Each leaf stores
$s_i=g(b_i)$. Each internal node $u$ stores
\[
  s_u=g(s_{u_L}\concat s_{u_R}),
\]
where $u_L$ and $u_R$ are its left and right children. At deployment, only
stored summaries need to move upward. For audit and validation, the raw span
beneath a node, written conceptually as $S(u)$, is retained or recoverable
from the original document.

\begin{figure}[t]
    \centering
    \includegraphics[width=0.92\linewidth]{01_base_plain.pdf}
    \caption{\textbf{A compression tree reduces a document to one root
    summary by recursively merging adjacent child summaries.} The v5 framing
    treats each stored summary as a learned sketch state. The root readout is
    valid only if the states preserve the information needed by $f$.}
    \label{fig:min-plain-tree}
\end{figure}

Write $u\sim_f v$ when the task readout cannot distinguish the two strings:
\[
  \metric\big(f(u),f(v)\big)=0.
\]
For approximate implementations, $\sim_f$ is replaced by a tolerance in the
task metric. Preservation is always relative to this task relation. A valid
summary keeps whatever evidence can move $f$ and may discard everything else.

\begin{assumption}[Context-compatible readout]\label{ass:context}
For every context $w\in\Strings$, $u\sim_f v$ implies both
$w\concat u\sim_f w\concat v$ and
$u\concat w\sim_f v\concat w$.
\end{assumption}

Context compatibility is an application assumption. It is plausible for
content rubrics whose evidence can be localized and substituted. It can fail
for targets such as tone, irony, or persuasion, where a locally equivalent
span may change how neighboring text is read.

\subsection{The Three Local Laws}

\paragraph{C1: Leaf sufficiency.}
\begin{equation*}\tag{C1}
  \forall\text{ realized leaf }b:\qquad g(b)\sim_f b.
\end{equation*}
A leaf summary must be sufficient for $f$ on its source span. In the manifesto
example, a leaf summary can compress prose aggressively, but it must preserve
the commitments and qualifications a policy coder would use if scoring that
span directly. In Count-Min, C1 is the fact that the item updates in a stream
segment are represented in the counter state used by the frequency readout.

\paragraph{C2: Re-summary stability.}
\begin{equation*}\tag{C2}
  \forall\text{ stored summary }s\in\operatorname{range}(g):\qquad
  g(s)\sim_f s.
\end{equation*}
Re-summarization of an already reduced state preserves the task value. This
is idempotence relative to the readout, not equality of strings. A second
summary may be shorter, cleaner, or differently worded, but it should not
change the policy evidence available to $f$. In a deterministic classical
sketch, the analogous operation is usually trivial because the state is
already in canonical form. In an LLM pipeline, it must be checked because
models can drift when asked to summarize their own summaries.

\paragraph{C3: Merge consistency.}
\begin{equation*}\tag{C3}
  \forall u,v\in\Strings:\qquad
  u\concat v
  \sim_f g(u\concat v)
  \sim_f g\!\left(g(u)\concat g(v)\right).
\end{equation*}
The disjoint route--summarize children, concatenate child states, and
summarize again--must preserve the same task value as summarizing the raw
union. This is the local mergeable-sketch law. In Count-Min, the corresponding
statement is $g(A\uplus B)=g(A)+g(B)$ before the frequency readout. In the
manifesto case, C3 asks whether the merge summary preserved interactions
between child spans: for example, whether a tax promise in one child and a
public-spending qualification in the other still appear together in the state
that reaches the parent.

\begin{figure}[t]
    \centering
    \includegraphics[width=0.94\linewidth]{05_full.pdf}
    \caption{\textbf{The local laws audit whether the learned state is
    sufficient along the realized tree.} C1 checks leaf encoding, C2 checks
    stability under re-summary, and C3 checks whether the split-then-merge
    route preserves the same readout as the raw union.}
    \label{fig:v5-local-laws}
\end{figure}

The laws are deliberately local. They do not require the researcher to prove
that the learned summarizer is an associative monoid homomorphism on all
possible inputs. They require evidence that the realized calls made by the
deployed tree preserved the task value. This is weaker than the classical
streaming setting, where arbitrary merge schedules and adversarial stream
orders force global associativity. It is also the right target for empirical
long-document measurement, where the tree is observed and auditable.

\subsection{What the Laws Certify}

The exact-law picture is simple. If every leaf state is sufficient for its raw
span, stored states do not drift under re-summary, and every merge state is
sufficient for the raw union of its children, then induction up the tree gives
root preservation: the root state has the same readout as the full document.
Approximate laws replace equality with a residual in the task metric. The
audit estimates that residual on the realized population of leaves and
internal nodes.

This is why C-TreePO is a learned mergeable sketch rather than a generic
summarization heuristic. The model is not merely asked to write good
summaries. It is asked to produce states that behave like sufficient
statistics for a specified readout under repeated merge operations. The
readout, the compression map, and the audit are therefore tied together.
Changing the target $f$ changes what information $g$ must keep. A state that
is sufficient for economic policy may be insufficient for immigration policy;
a state that preserves sentiment may be insufficient for fiscal-program
details. The sketch is task-relative.

\begin{proposition}[Strict and state-level mergeable limits]\label{prop:mergeable-reduction}
If the task value itself admits an associative merge, then oracle-preserving
compression reduces to a strict homomorphic readout. More generally, if a
bounded state carries the information needed by the readout and that state has
an associative merge, C-TreePO reduces to a classical mergeable summary on the
state even when the scalar answer lacks its own merge rule.
\end{proposition}

The exact nesting statement has four ingredients. First, there is a state
space $\mathcal{S}$ and a validity relation
$V(D,s)$ saying that state $s$ is a valid summary of dataset or span $D$.
Second, the C-Tree summarizer $g$ must implement the two state operations:
on raw leaves it behaves like the classical builder, $g(D)=\mathrm{build}(D)$,
and on child states it behaves like state merge,
$g(s_L\concat s_R)=\mathrm{merge}(s_L,s_R)$, up to the chosen serialization
of states as text. Third, validity must be closed under union:
if $V(D_L,s_L)$ and $V(D_R,s_R)$, then
$V(D_L\concat D_R,\mathrm{merge}(s_L,s_R))$. Fourth, the root readout $f$
must factor through valid states:
\[
  V(D,s) \quad\Longrightarrow\quad f(s)=F(D),
\]
or, in the approximate version,
$\metric(f(s),F(D))\le \varepsilon$. These are the precise roles of $g$ and
$f$: $g$ builds and merges state, while $f$ is only applied after the root
state has been produced.

Under those conditions, C1 is the leaf instance of build validity and readout
correctness, and C3 is the merge instance of validity/readout correctness for
the raw union of the two child spans. C2 is not part of Agarwal's state-merge
definition itself; it is the additional C-TreePO stability condition needed
when the system may re-summarize an already produced state, as LLM pipelines
often do. Classical canonical states usually satisfy this by construction.
The strict scalar route is a separate special case in which $F(D_L)$ and
$F(D_R)$ themselves determine $F(D_L\concat D_R)$; most sketches, including
distinct count and many quantile summaries, need the state-level route instead.

The proposition locates the classical case inside the learned one. Count-Min,
HyperLogLog, Bloom filters, and quantile sketches are not metaphors for
C-TreePO; they are exact special cases where the sufficient state and merge
law were designed analytically. C-TreePO keeps the same target--state--readout
shape but permits $g$ to be learned in a richer space and audited after
deployment. Classical full mergeable summaries nest through this relational
state-level shape; randomized mergeable summaries nest in probability, because
root readout is correct on the event that the merged root state is valid for
the union stream. The approximation parameter $\varepsilon$ and any
training-data/sample-size requirement remain separate profile or concentration
conditions, not scalar answer-merge laws. For LLMs, the learned state is also
inspectable text. That
observability turns a black-box compression problem into a measurable one:
node labels, direct reading, and counterfactual fiber tests can all ask
whether $g(x)$ really kept what $f$ needs.

In the neural-operator route, $f$ is usually a calibrated readout
$\widehat f$ rather than the true oracle itself. The theorem then separates
readout calibration from state learning: learn $\widehat f$ so that it is
uniformly within $\varepsilon_f$ of $\fstar$, and learn $g$ by projecting leaf
and merge calls toward the state-level C1/C3/C2 laws measured by
$\widehat f$. If the realized $g$ is exact or approximately theorem-backed
for $\widehat f$, the true-oracle distortion for $\fstar$ is the local-law
budget plus the additive calibration slack $2\varepsilon_f$. Neural
operators, including transformer/LLM-style operators at the architecture
level, supply the ambient function class; they do not by themselves certify
mergeability or scalar child-answer merging.
